% ============================================================
% chapters/05-alternatives.tex -- פרק 5: חלופות שיטה, בחירה וחדשנות
% Content only. All design comes from template.cls.
% Figures: figures/ch5-architecture.tex (TikZ),
%          figures/ch5-decision.pdf, figures/ch5-sensitivity.pdf (code/make_figures.py)
% The scores and weights in tab:decision MUST match ALTS and WEIGHTS in
% code/make_figures.py. Change both together.
% ============================================================

\chapter{חלופות שיטה, בחירת הפתרון וחדשנות}
\label{ch:alternatives}

בפרויקט מחקרי, ״חלופות התכן״ הן דרכים שונות לענות על שאלת המחקר. פרק זה מציג שלוש חלופות ששונות זו מזו בעיקרון ולא רק בפרטים, בוחר ביניהן במטריצת החלטה משוקללת, בודק אם הבחירה יציבה, ומתאר את הארכיטקטורה של החלופה שנבחרה ואת החדשנות שבה.

\section{שלוש חלופות שונות עקרונית}
\label{sec:alt-three}

\paragraph{חלופה א: מחקר מקורות קלאסי.} זו הדרך שבה נחקרה השאלה עד היום. החוקר בוחר מקורות מרכזיים מכל אחת משלוש המדינות – כרוניקות, מכתבים, חוקים והחלטות של מועצות כנסייה – קורא אותם לעומק, ומשווה ביניהם בכתיבה פרשנית. התוצר הוא חיבור טיעוני, בדומה לעבודות שנסקרו בפרק~\ref{ch:literature}.

\paragraph{חלופה ב: מאגר נתונים וקידוד ידני.} צוות של מומחים בונה מאגר נתונים לפי המשתנים שהוגדרו בפרק~\ref{ch:problem}. כל ערך במאגר נקבע בידי מומחה שקורא את המקורות הרלוונטיים ומקודד אותם לפי ספר הנחיות אחיד. זו הגישה של מאגר \en{Seshat} ושל המחקרים שמדדו את התנצרות העילית הרומית \cite{salzman2002making}. התוצר הוא מאגר מובנה, שעליו מריצים ניתוח סטטיסטי.

\paragraph{חלופה ג: מאגר נתונים וניתוח טקסט בעזרת \en{AI}, עם תיקוף מומחה.} כמו בחלופה ב, נבנה מאגר מובנה; אבל את מדדי ההשתלבות בין הכנסייה למדינה מפיק מודל שפה, שמסווג אלפי קטעים מתוך המקורות. מומחה בודק מדגם מתוך הסיווגים, ורק אם רמת ההתאמה עומדת ביעד (דרישה \en{REQ-AI-01}), התוצאות נכנסות למאגר. משתני הפיקוח – כלכלה, צבא ואקלים – נלקחים מהספרות ומ־\en{Seshat}.

טבלה~\ref{tab:alt-tradeoffs} מסכמת את היתרון ואת המחיר של כל חלופה.

\begin{table}[htbp]
  \centering
  \caption{שלוש החלופות: היתרון והמחיר של כל אחת}
  \label{tab:alt-tradeoffs}
  \small
  \begin{tabular}{p{30mm}p{55mm}p{55mm}}
    \toprule
    \textbf{חלופה} & \textbf{היתרון} & \textbf{המחיר} \\
    \midrule
    א: מחקר מקורות קלאסי & עומק פרשני, הבנת הקשר, עלות כספית נמוכה & כיסוי חלקי, דוגמאות נבחרות, אין בדיקה סטטיסטית של ההשערות \\
    \midrule
    ב: מאגר וקידוד ידני & דיוק גבוה בכל ערך, שיטה מוכרת ומקובלת & אלפי שעות עבודה של מומחים; לכן אפשר לקודד רק חלק מהטקסטים \\
    \midrule
    ג: מאגר וניתוח \en{AI} עם תיקוף & כיסוי של אלפי טקסטים באותה שיטה, שחזור מלא של התהליך & סיכון לטעויות של המודל, עלות טוקנים, צורך בתיקוף \\
    \bottomrule
  \end{tabular}
\end{table}

\section{מטריצת החלטה משוקללת}
\label{sec:alt-matrix}

\begin{keybox}{קודם המשקלים, אחר כך הציונים}
מטריצה שבה המשקלים נקבעים אחרי שכבר ידוע מי צריך לנצח היא קישוט ולא כלי החלטה. לכן הקריטריונים והמשקלים נקבעו לפני מתן הציונים, ובהמשך נבדק אם ההחלטה נשארת יציבה כשהמשקלים משתנים (סעיף~\ref{sec:alt-sensitivity}).
\end{keybox}

\subsection*{הקריטריונים והמשקלים}

\begin{itemize}
  \item \textbf{תוקף מדעי (0.30):} עד כמה התוצאות משקפות נכון את המקורות, ומאפשרות לאשש או להפריך את ההשערות. זה הקריטריון החשוב ביותר במחקר, ולכן קיבל את המשקל הגבוה ביותר.
  \item \textbf{יכולת שחזור (0.20):} האם חוקר אחר יכול לחזור על התהליך ולקבל אותן תוצאות.
  \item \textbf{כיסוי המקורות (0.20):} איזה חלק מהטקסטים הרלוונטיים בשלוש המדינות נכנס לניתוח.
  \item \textbf{עלות (0.15):} שעות עבודה של חוקרים ועלות טוקנים. ציון גבוה פירושו עלות נמוכה.
  \item \textbf{חדשנות (0.15):} מה החלופה מוסיפה על מה שכבר נעשה בספרות.
\end{itemize}

\subsection*{הציונים}

טבלה~\ref{tab:decision} מציגה את הציון של כל חלופה בכל קריטריון, בסולם של 1 עד 5, ובסוגריים את הציון המשוקלל.

\begin{table}[htbp]
  \centering
  \caption{מטריצת החלטה משוקללת: ציון (ובסוגריים הציון המשוקלל)}
  \label{tab:decision}
  \small
  \begin{tabular}{p{38mm}ccc}
    \toprule
    \textbf{קריטריון (משקל)} & \textbf{א: מקורות} & \textbf{ב: קידוד ידני} & \textbf{ג: \en{AI} ותיקוף} \\
    \midrule
    תוקף מדעי (0.30)    & 3 (0.90) & 5 (1.50) & 4 (1.20) \\
    יכולת שחזור (0.20)  & 2 (0.40) & 4 (0.80) & 5 (1.00) \\
    כיסוי המקורות (0.20) & 2 (0.40) & 3 (0.60) & 5 (1.00) \\
    עלות (0.15)         & 5 (0.75) & 2 (0.30) & 3 (0.45) \\
    חדשנות (0.15)       & 1 (0.15) & 3 (0.45) & 5 (0.75) \\
    \midrule
    \textbf{סה״כ}       & \textbf{2.60} & \textbf{3.65} & \textbf{4.40} \\
    \bottomrule
  \end{tabular}
\end{table}

\paragraph{נימוק הציונים.} בתוקף המדעי, הקידוד הידני (ב) קיבל את הציון הגבוה ביותר: בהנחות התכנון מיוחס למומחה שקורא כל קטע יתרון בתוקף על פני סיווג אוטומטי; זו הנחה שתיבדק בפיילוט. חלופה ג קיבלה ציון נמוך ממנה בנקודה אחת, משום שגם אחרי תיקוף נשאר סיכון לטעויות המודל; חלופה א קיבלה 3, משום שהיא אינה בודקת את ההשערות סטטיסטית. ביכולת שחזור ובכיסוי, חלופה ג חזקה במיוחד: אותו תהליך רץ על כל הטקסטים באותו אופן, ואפשר להריץ אותו שוב. בעלות, המחקר הקלאסי (א) הוא הזול ביותר, והקידוד הידני (ב) היקר ביותר בגלל אלפי שעות של מומחים; חלופה ג נמצאת באמצע, משום שהיא מחליפה חלק משעות העבודה בעלות של טוקנים. בחדשנות, חלופה א אינה מוסיפה שיטה חדשה, וחלופה ג מציעה לשלב את שני הכלים לצורך שאלת המחקר.

איור~\ref{fig:alt-decision} מפרק את הציון הסופי של כל חלופה לתרומה של כל קריטריון. אפשר לראות שהיתרון של חלופה ג נבנה בעיקר מכיסוי, משחזור ומחדשנות, ואילו בתוקף המדעי היא מפסידה לחלופה ב.

\begin{figure}[htbp]
  \centering
  \includegraphics[width=0.85\textwidth]{figures/ch5-decision.pdf}
  \caption{הציון הסופי של כל חלופה, מפורק לתרומת כל קריטריון}
  \label{fig:alt-decision}
\end{figure}

\section{ניתוח רגישות: האם ההחלטה יציבה?}
\label{sec:alt-sensitivity}

רגישות ההחלטה למשקלים נבדקה בשתי דרכים, באמצעות הקוד המצורף.

\paragraph{שינוי של קריטריון אחד.} נבחר לבדיקה את התוקף המדעי, משום שזה הקריטריון שבו חלופה ג חלשה מחלופה ב. משקלו הועלה בהדרגה מ־0 עד 0.9, ושאר המשקלים הוקטנו באופן יחסי. כפי שמראה איור~\ref{fig:alt-sensitivity}, חלופות ב ו־ג משתוות כאשר משקל התוקף המדעי הוא 0.60; ברשת החישוב, בצעד של 0.01, חלופה ג מפסיקה להוביל ב־0.61 – יותר מפי שניים ממשקל הבסיס. רק כשהתוקף המדעי שווה יותר מכל שאר הקריטריונים יחד, הקידוד הידני עדיף.

\begin{figure}[htbp]
  \centering
  \includegraphics[width=0.85\textwidth]{figures/ch5-sensitivity.pdf}
  \caption{ניתוח רגישות: הציון של כל חלופה כשמשקל התוקף המדעי משתנה}
  \label{fig:alt-sensitivity}
\end{figure}

\paragraph{משקלים אקראיים.} בבדיקה השנייה הוגרלו 10,000 צירופים של משקלים סביב משקלי הבסיס, וחושב בכל פעם איזו חלופה מנצחת. ב־10,000 הגרלות עם זרע קבוע 1, חלופה ג ניצחה ב־9,960 מקרים (99.60\%), חלופה ב ב־38 (0.38\%), וחלופה א בשניים (0.02\%). המשקלים הוגרלו מהתפלגות \en{Dirichlet} עם פרמטרים השווים למשקלי הבסיס כפול 20. שיעור זה מתאר רגישות להנחות הציונים והמשקלים, ואינו הסתברות שהמחקר יצליח. המסקנה היא שההחלטה יציבה: היא אינה תוצאה של בחירה מקרית של משקלים.

\paragraph{מה הבדיקה מלמדת.} מעבר לבחירה עצמה, הניתוח מצביע על נקודת התורפה של החלופה שנבחרה: התוקף המדעי של סיווגי המודל. לכן הארכיטקטורה (סעיף~\ref{sec:alt-architecture}) בנויה סביב תיקוף אנושי, ולכן הטעויות של המודל הן אחד הסיכונים המרכזיים בפרק~\ref{ch:risks}.

\section{ארכיטקטורת המחקר}
\label{sec:alt-architecture}

איור~\ref{fig:alt-architecture} מציג את מסלול הנתונים בחלופה שנבחרה, מהמקורות ועד לתוצאות.

\begin{figure}[htbp]
  \centering
  \input{figures/ch5-architecture}
  \caption{ארכיטקטורת המחקר: מסלול הנתונים מהמקורות ועד לתוצאות}
  \label{fig:alt-architecture}
\end{figure}

\begin{enumerate}
  \item \textbf{מקורות.} הטקסטים נאספים מ־\en{Perseus} ומ־\en{dMGH}, ששניהם בגישה פתוחה ובפורמט שהמחשב יכול לקרוא (פרק~\ref{ch:market}).
  \item \textbf{הכנת הקורפוס.} כל טקסט מחולק לקטעים קצרים, וכל קטע מקבל תגיות: מקור, מדינה, ומקטע הזמן של 25 שנה שאליו הוא שייך.
  \item \textbf{סיווג בעזרת מודל שפה.} המודל מסווג כל קטע לפי קטגוריות שנגזרות מהמשתנים בטבלה~\ref{tab:variables}: למשל, בישוף בתפקיד ממשלתי, חקיקה דתית מטעם השליט, או סכסוך בין שליט לבישופים.
  \item \textbf{תיקוף בידי מומחה.} מומחה לתקופה יתייג מדגם פיתוח נפרד, ולאחר הקפאת ההנחיות יתייג מדגם מבחן של 300 קטעים בלי לראות את תשובות המודל. מדגם המבחן לא ישמש לתיקון ההנחיות. אם ההתאמה נמוכה מ־85\%, הסיווג אינו מאושר; תיקון יחייב מבחן חדש או מעבר לקידוד ידני. פרטי ההפרדה מפורטים בסעיף~\ref{sec:alt-validation}.
  \item \textbf{מאגר הנתונים.} הסיווגים המאושרים נקשרים לרשומת אירוע מאוחדת לפני האיגום האזורי; ריבוי אזכורים לא ייחשב כריבוי אירועים. לצד האירועים נכנסים משתני הפיקוח.
  \item \textbf{ניתוח סטטיסטי.} על המאגר נבדקות ההשערות \en{H1} ו־\en{H2}, עם ניתוח רגישות להגדרות השונות של המדדים.
  \item \textbf{תוצאות ופרסום.} הקוד, הנתונים והקורפוס המתויג מתפרסמים בגישה פתוחה.
\end{enumerate}

\subsection*{ועדת הסוכנים}

שלב הסיווג אינו נעשה בידי מודל אחד, אלא בידי צוות קטן של סוכני \en{AI}, שכל אחד מהם מקבל תפקיד מוגדר בקובץ הנחיות משלו. \textbf{הסוכן המסווג} קורא כל קטע ומסווג אותו. \textbf{הסוכן המבקר}, בתפקיד ״צוות אדום״, מקבל מדגם מהסיווגים ומחפש טעויות, סתירות וסיווגים חשודים. \textbf{הסוכן הבודק} מוודא שהנתונים שלמים ועקביים: שלכל קטע יש מקור ותאריך, ושאין כפילויות. סוכן ראשי מחלק את העבודה ביניהם ומרכז את התשובות. הכלל המנחה הוא הפרדה בין מי שמגלה תקלה לבין מי שמתקן אותה, וההחלטה הסופית נשארת תמיד בידי המומחה האנושי.

\section{החדשנות}
\label{sec:alt-innovation}

המחקר המוצע אינו שכפול של עבודה קיימת. החדשנות שלו נמצאת בארבעה מקומות:

\begin{enumerate}
  \item \textbf{השוואה של שלושה מקרים באותם מדדים.} לפי הסקירה בפרק~\ref{ch:literature}, המחקר מציע לבחון את רומא המערבית, רומא המזרחית והממלכה הפרנקית בשיטה כמותית משותפת. הסקירה מצביעה על פער רלוונטי, אך אינה מוכיחה שלא קיימת עבודה קודמת בעלת חפיפה.
  \item \textbf{מדידה של השתלבות, ולא רק של התנצרות.} מחקרים קודמים מדדו כמה אנשים התנצרו; המחקר המוצע מודד \textbf{איך} הכנסייה השתלבה במוסדות המדינה.
  \item \textbf{שילוב של מודל שפה ומומחה.} המחקר מאמץ את המסקנה של מחקרים כמו \en{Ithaca}, שלפיה אדם ומודל יחד עדיפים על כל אחד מהם לבד \cite{assael2022ithaca}, ומיישם אותה בשאלה היסטורית גדולה.
  \item \textbf{תוצרים פתוחים.} המאגר והקורפוס המתויג יעמדו לרשות חוקרים אחרים, ויאפשרו לבדוק גם השערות אחרות על אותה תקופה.
\end{enumerate}


\section{תוכנית תיקוף ובקרת שגיאות}
\label{sec:alt-validation}

הנחיות הקידוד יוגדרו לפני מבחן הסיווג ויכללו דוגמאות חיוביות, דוגמאות נגדיות ואפשרות ״לא ניתן לקבוע״. לשלב הפיתוח יוקצו 150 קטעים, 50 מכל מדינה. לאחריו יישמרו 300 קטעי מבחן נפרדים, 100 מכל מדינה, תוך פיזור לפי תקופה וסוג מקור. החלוקה תיעשה בקבוצות לפי החיבור המקורי: כל המהדורות, התרגומים והקטעים המקבילים של אותו חיבור יישארו באותה קבוצה. בדיקת דמיון תאתר גם קטעים סמוכים או כמעט זהים לפני החלוקה. אם פיזור זה אינו אפשרי, הפער ידווח לפני המבחן ולא יוסתר בציון הכולל.

המומחה יתייג את קטעי המבחן בלי לצפות בפלט המודל. מתוך המבחן, 60 קטעים יתויגו גם בידי מומחה שני; מחלוקות יתועדו ויוכרעו לפי ספר הקידוד לפני חשיפת תשובות המודל. ידווחו התאמה כוללת, מטריצת שגיאות ותוצאות נפרדות לכל מדינה ולכל קטגוריה. קטגוריה נדירה לא תיחשב מוצלחת רק משום שהמודל מסווג את כל הדוגמאות כשליליות. ערכי אי־ודאות יישמרו לצד התיוגים, ולא יומרו אוטומטית לאפס.

סף הקבלה הוא התאמה מלאה של וקטור התוויות ב־85\% לפחות מ־300 קטעי המבחן: קטע נספר כנכון רק אם כל התוויות שנקבעו עבורו תואמות לתיוג המומחה. ״לא ניתן לקבוע״ היא תווית מפורשת, והסכמה עליה תדווח גם בנפרד; אין להשמיט קטעים קשים כדי להעלות את הציון. המדדים לכל תווית הם אבחוניים ואינם מחליפים את מדד הקבלה הראשי, בהתאם לדרישה \en{REQ-AI-01}. בדיקת העקביות תיערך על אותם קטעים בשתי הרצות עם גרסת מודל והנחיות קבועות, תוך מדידת שיעור שינוי התוויות, בהתאם ל־\en{REQ-AI-02}. לאחר כישלון לא יבוצע כוונון על מדגם המבחן שכבר נחשף. יוקצה מדגם חדש של 300 קטעים או שיופעל מסלול הקידוד הידני; התוצאה המקורית תישמר בדוח הבדיקה.

\section{תוכנית ניתוח וגבולות ההסקה}
\label{sec:alt-statistics}

יחידת הניתוח היא מדינה במקטע זמן, ולא קטע טקסט. אלפי קטעים אינם אלפי תצפיות עצמאיות על יציבות המדינה. במאגר מתוכננים כ־34 מקטעים בלבד; שלוש המדינות אינן מדגם אקראי, והפרנקים פועלים בחלקם בתקופה מאוחרת משני המקרים הרומיים. לכן מטרת הניתוח היא לבחון קשרים התואמים להשערות, תוך הצגת הסברים חלופיים. אין להסיק ממנו לבדו השפעה סיבתית של הנצרות.

לפני בדיקת ההשערות יוגדרו נוסחת מדד ההשתלבות, כללי הספירה, אורך החשיפה ומדיניות הערכים החסרים. המדד התלוי הראשי יהיה מספר ניסיונות תפיסת שלטון ומרידות ביחס למספר השנים המכוסות. יוצגו תחילה סדרות זמן ונתונים נפרדים לכל מדינה. אם הכיסוי והשונות מאפשרים זאת, יושוו שני מודלי ספירה חקרניים ממשפחת \en{Poisson}: אחד עם מדד התנצרות לצורך בחינת \en{H1}, והשני עם מדד השתלבות לצורך בחינת \en{H2}. בשניהם יופיעו אותו משתנה פיקוח שנקבע מראש, אותם אינדיקטורים אזוריים ואותו אורך חשיפה. במודל הספירה לוגריתם שנות הכיסוי ישמש \en{offset}; מקטע חלקי לא יטופל כמקטע מלא. כאשר מדד ההתנצרות חסר או כמעט קבוע, לא תוצג \en{H1} כהשערה שנבחנה בהצלחה. לא יוכנסו כל 12 המשתנים בו־זמנית למודל זעיר. ייבדקו פיזור יתר ותלות בזמן; אם הנחות המודל אינן סבירות, לא ידווחו ערכי מובהקות כאילו היו תקפים.

רמת הבדיקה המתוכננת היא 0.05, אך הצלחת המחקר אינה מותנית בתוצאה מובהקת. ידווחו גודל האפקט, אי־הוודאות, מספר המקטעים שנכללו ושינויים בתוצאה כאשר מושמטים מקורות או משתנות הגדרות המדדים. חוסר מידע אינו היעדר אירוע: לכל ערך חסר יירשמו סיבה ומקור. איחוד מקטעים ל־50 שנה יוצג כניתוח רגישות נוסף לצד טבלת 25 השנים; הוא אינו מבטל את דרישת הכיסוי המקורית או הופך נתון חסר לנתון ידוע.


\paragraph{איגום אירועים ומקורות.} מאגר אירועים יחיד יקצה לכל מרידה, מועצה או החלטה מזהה קבוע, תאריך וטווח אי־ודאות, תחולה טריטוריאלית ורשימת כל מראי המקום התומכים בו. כמה כרוניקות המתארות אותה מרידה אינן כמה מרידות. לכל צירוף אירוע–אזור יישמר סימון אם מדובר באירוע אזורי או משותף; שורות משותפות לא ישמשו הוכחה לשתי תצפיות עצמאיות. לצורך \en{REQ-DATA-02} תיערך בדיקה אוטומטית של נוכחות מזהה מקור ומראה מקום בכל ערך, לצד ביקורת מדגמית של התאמת התוכן. רשומה חסרת תיעוד תסומן כלא מאומתת ולא תיכנס לנתונים המאושרים.
